\Needspace{16\baselineskip}
\Article{356}
\textsc{Satz.} \emph{Es sei $P$ eine Primfunktion von $x$ vom Grad $m$ und $X$ eine Funktion von $x,x^p,x^{pp},x^{p^3}\dotstwo x^{p^{m-1}}$, in der alle diese Größen auf gleiche Weise vorkommen, d. h. die unverändert bleibt, wie auch immer diese Größen untereinander vertauscht werden. Dann ergibt die Funktion $X$ bei der Division durch $P$ einen Rest, der eine Zahl ist.}

\emph{Beweis.} Es sei der Rest
\[
Ax^{m-1}+Bx^{m-2}+\dotstwo+N\equiv\xi
\]
Alle Koeffizienten $A,B,C\ldots$ bis ausschließlich $N$ sind $\equiv0$. Dies wird folgendermaßen bewiesen. Da $X-\xi$ durch $P$ teilbar ist, lässt sich auch $X^p-\xi^p$ durch $P$ teilen.
\SourceLeaf{238}
Man erkennt aber leicht, dass $X^p$ aus $X$ entsteht, wenn man $x$ durch $x^p$, $x^p$ durch $x^{pp}$ usw.\ .\ und $x^{p^{m-1}}$ durch $x^{p^m}$ oder, was dasselbe ist, durch $x$ ersetzt. Somit ist offenbar $X^p\equiv X\;(\mathrm{mod.}\,P)$. Da nun $X^p\equiv\xi^p$ und $X\equiv\xi(\mathrm{mod.}\,P)$, ist auch $\xi^p\equiv\xi(\mathrm{mod.}\,P$, oder
\[
\xi^p-\xi\equiv0(\mathrm{mod.}\,P)
\]
Doch ist $\xi^p-\xi$ nach dem Modul $p$ dem Produkt aus $\xi,\xi+1,\xi+2,\dotstwo$ bis $\xi+p-1$ kongruent. Diese Faktoren sind alle zu $P$ teilerfremd, sofern $\xi$ nicht einfach eine Zahl ist. Daher kann auch $\xi^p-\xi$ auf keine andere Weise durch $P$ teilbar sein. Q. E. D.

Solche Funktionen sind die Summe aller Größen, die Summe ihrer Quadrate, Kuben usw. und die Summen der Produkte aus je zwei, drei usw. Größen. Welche Zahl dies aber ist, werden wir im folgenden Artikel bestimmen:

\Needspace{16\baselineskip}
\Article{357}
\textsc{Satz.} \emph{Die Primfunktion des vorigen Artikels sei}
\[
P\equiv x^m-Ax^{m-1}+Bx^{m-2}-Cx^{m-3}+\text{usw.}
\]
\emph{Dann ist der Rest bei der Division der Summe der Größen $x,x^p$ usw. $x^{p^{m-1}}$ durch $P$ $\equiv A$, bei der Summe der Produkte aus je zwei $\equiv B$, bei der Summe der Produkte aus je drei $\equiv C$ usw.}

\emph{Beweis.} Jene Funktionen seien $X,Y,Z$ usw. und ihre Reste der Reihe nach die Zahlen $A',B',C'$ usw. Man sieht leicht, dass $x,x^p,x^{pp}$ usw. Wurzeln der Gleichung
\[
z^m-Xz^{m-1}+Yz^{m-2}-Zz^{m-3}+\text{usw.}=0
\]
sind. Setzt man daher $z=x$, so ist
\[
x^m-Xx^{m-1}+Yx^{m-2}-Zx^{m-3}+\text{usw.}=0
\]
Die Funktionen $X-A'$, $Y-B'$, $Z-C'$ usw. sind aber durch $P$ teilbar und daher auch die Funktion
\[
x^m-A'x^{m-1}+B'x^{m-2}-C'x^{m-3}+\text{usw.}
\]
Dies ist nur möglich, wenn $A'\equiv A$, $B'\equiv B$, $C'\equiv C$ usw. Q. E. D.

Ferner ist bekannt, dass jede andere Funktion $X$ von $x,x^p$ usw. [in der alle diese Größen auf gleiche Weise vorkommen,] stets aus diesen abgeleitet werden kann. So ist
\SourceLeaf{239}
\[
x^2+x^{2p}+x^{3pp}+\text{usw.}\equiv AA-2B(\mathrm{mod.}\,P)\text{ usw. usw.}
\]
Beispiel. Es sei $p=5$ und $P\equiv x^2+2x+3$. Dann ergibt die Funktion $x+x^5$ bei der Division durch $P$ den Rest $\equiv-2$, und es ist $x^6\equiv3$ usw. usw.
